\section{Further coordinate block estimates and spectral stability}
\label{sec:additional-results}

The estimates in this section apply to blocks of original independent
coordinates with densities bounded by one.  Their conclusions are
uniform in affine terms.  The deterministic spectral inequalities
control perturbations of coordinate submatrices.

\subsection{A linear bank with two crossing directions}

\begin{lemma}[Negative moments from a density bound]
\label{supp:negative-moment}
If a probability law on $\mathbb R^d$ has density at most $K$, then for
$0<p<d$ and every $z\in\mathbb R^d$,
\[
 \mathbb E|Y-z|^{-p}\le\frac d{d-p}(K v_d)^{p/d},
\]
where $v_d$ is the volume of the unit ball.
\end{lemma}
\begin{proof}
For every $r>0$, the probability of the ball of radius $r$ is at most
$\min\{1,Kv_dr^d\}$.  Integrate this bound in
\[
 \mathbb E|Y-z|^{-p}
 =\int_0^\infty\mathbb P\{|Y-z|<t^{-1/p}\}\,dt
\]
and split at $t=(Kv_d)^{p/d}$.  The two contributions are
$(Kv_d)^{p/d}$ and $p(Kv_d)^{p/d}/(d-p)$.
\end{proof}

\begin{theorem}[An arbitrary-rank crossing block with a linear bank]
\label{supp:null-bank}
Let $X_I$ and $X_J$ be disjoint original coordinate banks and let
\[
 Q=F(X_I)+2X_I^{\mathsf T}M X_J+b_J^{\mathsf T}X_J,
\]
where $F$ is any real measurable function.  If $s_2(M)>0$, then
\begin{equation}\label{eq:supp-null-bank}
 \|p_Q\|_\infty\le
 \frac{2\sqrt\pi}{\sqrt{s_1(M)s_2(M)}}.
\end{equation}
The rank and all smaller singular values of $M$ are unrestricted.
\end{theorem}
\begin{proof}
Condition on $X_I=x$.  The remaining polynomial is linear in $X_J$,
with coefficient vector $c(x)=2M^{\mathsf T}x+b_J$.  The
one-dimensional product-projection estimate of \citet{LPP2016} gives
a conditional density bound $\sqrt2/|c(x)|$ whenever $c(x)\ne0$.
Let $u_1,u_2$ and $v_1,v_2$ be corresponding left and right singular
vectors for the two largest singular values $s_1,s_2$ of $M$.
The vector $U=(u_1^{\mathsf T}X_I,u_2^{\mathsf T}X_I)$ has density
bounded by two.  Therefore
\[
 Z=(2s_1U_1+v_1^{\mathsf T}b_J,
       2s_2U_2+v_2^{\mathsf T}b_J)
\]
has density at most $1/(2s_1s_2)$, and $|c(X_I)|\ge|Z|$.
Lemma~\ref{supp:negative-moment} with $d=2,p=1$ gives
\[
 \mathbb E|c(X_I)|^{-1}
 \le2\sqrt{\frac\pi{2s_1s_2}}.
\]
In particular, the exceptional set $c(X_I)=0$ has probability zero.
Averaging the conditional density inequality proves
\eqref{eq:supp-null-bank}.  This argument retains the full matrix $M$;
the projection onto two singular directions is used only for a lower
bound on the length of $c(X_I)$.
\end{proof}

The condition that the $J$ bank enters linearly is essential to this
proof.  Theorem~\ref{supp:null-bank} includes a bilinear form with
an arbitrary quadratic or measurable addition on one bank.  Its
conditioning argument is closely related to the independent-block
density estimates in \citet{BGKS2024}.

\subsection{A singular physical pair and a random slope}

\begin{theorem}[Rank-one physical pair]
\label{supp:rank-one-pair}
Let $X,Y,Z_1,\ldots,Z_m$ be independent original coordinates.  Suppose
\begin{align*}
 Q={}&aX^2+2bXY+cY^2+2X\langle u,Z\rangle
       +2Y\langle v,Z\rangle+F(Z)\\
    &+2\alpha X+2\beta Y,
\end{align*}
where $a\ne0$, $ac=b^2$, $F$ is real measurable, and
$w=v-(b/a)u\ne0$.  Then
\begin{equation}\label{eq:supp-rank-one-pair}
 \|p_Q\|_\infty\le
 \frac{4\,2^{1/4}}{\sqrt{|a|\,|w|}}.
\end{equation}
The dependence of $F$ on all coordinates of $Z$ is unrestricted.
\end{theorem}
\begin{proof}
Complete the square only in the original coordinate $X$.  With
$\eta=\beta-b\alpha/a$ this writes
\[
 Q=a\left(X+\frac{bY+\langle u,Z\rangle+\alpha}{a}\right)^2
       +2Y(\langle w,Z\rangle+\eta)+\widetilde F(Z).
\]
Conditioning on $(Y,Z)$ and using the scalar square density bound gives
\[
 p_Q(t)\le |a|^{-1/2}
  \mathbb E k_a\bigl(t-2Y(\langle w,Z\rangle+\eta)
                     -\widetilde F(Z)\bigr),
 \quad k_a(u)=|u|^{-1/2}\mathbf1_{\{au>0\}},
\]
for almost every $t$.  For any probability density $f\le M$,
\[
 \sup_t\int_0^\infty y^{-1/2}f(t\pm y)\,dy\le2\sqrt M,
\]
by filling the interval of length $1/M$ nearest zero.
If $\ell\ne0$, the variable $2\ell Y$ has density at most
$1/(2|\ell|)$.  Hence the expectation of the one-sided kernel
over $Y$ is at most $\sqrt2\,|\ell|^{-1/2}$, uniformly in its
translate.
The scalar projection $\langle w,Z\rangle+\eta$ has density at
most $\sqrt2/|w|$ by \citet{LPP2016}.
Lemma~\ref{supp:negative-moment} with $d=1,p=1/2$ gives
\[
 \mathbb E|\langle w,Z\rangle+\eta|^{-1/2}
 \le2\sqrt{2\sqrt2/|w|}.
\]
Multiplication by the preceding bound for the slope gives
$\sqrt2\cdot2\sqrt{2\sqrt2}=4\,2^{1/4}$, as asserted.
All changes of variables can be integrated against nonnegative output
test functions, so the almost-everywhere conditional formulas introduce
no exceptional-level assumption.
\end{proof}

\begin{corollary}[An aggregate bound for singular pairs]
\label{supp:singular-pair-volume}
Under Theorem~\ref{supp:rank-one-pair}, put
\[
 W=a^2\sum_{i\ne j}w_i^2w_j^2.
\]
If $W>0$, then
\[
 \|p_Q\|_\infty\le
     4\,2^{1/4}|a|^{-1/4}W^{-1/8}.
\]
\end{corollary}
\begin{proof}
The inequality $|w|^4\ge\sum_{i\ne j}w_i^2w_j^2=W/a^2$
inserted into \eqref{eq:supp-rank-one-pair} proves the statement.
\end{proof}

\subsection{Spectral energy in a coordinate set}

For a coordinate set $H$, let $P_H:\mathbb R^{|H|}\to\mathbb R^n$
be the coordinate inclusion.  Compression is written
$T_{HH}=P_H^{\mathsf T}TP_H$.  Eigenvalues of a positive
semidefinite matrix are listed in decreasing order.

\begin{lemma}[Spectral-square spill]
\label{supp:spill}
Let $A$ be symmetric, $C=A_{HH}$, $B=A_{H,H^c}$, and
$0\le G\le A^2$.  Suppose $\lambda_k(G_{HH})\ge g^2$.
For $0\le r<k$ put $c=s_{r+1}(C)$.  If $c<g$, then
\[
 s_{k-r}(B)\ge\sqrt{g^2-c^2}.
\]
In particular, one may use $G=\Lambda^2
\mathbf1_{\{|A|\ge\Lambda\}}$ for any $\Lambda>0$.
\end{lemma}
\begin{proof}
The exact block identity is $(A^2)_{HH}=C^2+BB^{\mathsf T}$.
Write $C=C_r+E$, where $C_r$ consists of its $r$ largest
absolute eigenmodes.  Then $C_rE=0$, $\operatorname{rank}C_r\le r$,
and $E^2\le c^2I$.  Thus
\[
 BB^{\mathsf T}+C_r^2=(A^2)_{HH}-E^2\ge G_{HH}-c^2I.
\]
The left side has at least $k$ eigenvalues at least $g^2-c^2$.
Adding a positive semidefinite matrix of rank at most $r$ can
increase that count by at most $r$.  Consequently
$BB^{\mathsf T}$ has at least $k-r$ eigenvalues of that size,
which is the assertion.  The stated spectral projection gives an
admissible $G$ by the spectral theorem.
\end{proof}

\begin{lemma}[A pair in an approximate rank-two self block]
\label{supp:pair-extraction}
Let $C=C^{\mathsf T}$ be $m$ by $m$, $m\ge3$, and write
$N=\binom m2$.  If $s_2(C)\ge b$ and $s_3(C)<a$, then some
principal two-coordinate block $K$ satisfies
\[
 s_{\min}(K)>b/N-a.
\]
\end{lemma}
\begin{proof}
Let $C_2=\lambda_1uu^{\mathsf T}+\lambda_2vv^{\mathsf T}$ be the
two largest absolute eigenmodes.  The orthonormality of $u,v$ gives
$\sum_{i<j}(u_iv_j-u_jv_i)^2=1$.  Choose a pair $I$ whose
summand is at least $1/N$.  Then
\[
 |\det(C_2)_{II}|\ge s_1(C)s_2(C)/N,
 \qquad \|(C_2)_{II}\|\le s_1(C).
\]
It follows that $s_{\min}((C_2)_{II})\ge b/N$.
The remaining matrix has norm $s_3(C)<a$, so the singular-value
perturbation inequality proves the claim.
\end{proof}

\begin{theorem}[Six spectral-square directions on a coordinate set]
\label{supp:six-carrier}
Let $s_3(A)\ge s>0$ and let $H$ have size $m\ge6$.  Suppose
$0\le G\le A^2$ and $\lambda_6(G_{HH})\ge g^2$.
For $N=\binom m2$, choose $0<a<b/N$ and $0<b<g$, and set
$\kappa=b/N-a$.  Then
\begin{equation}\label{eq:supp-six-carrier}
 \|p_Q\|_\infty\le\max\left\{
 \frac{2\sqrt{6\pi m}}a,
 \frac{9.2+\log_+(10000\kappa/s)}\kappa,
 \frac{5\sqrt{2e}}{\sqrt{g^2-b^2}}\right\}.
\end{equation}
All entries of the actual matrix $A$ outside $H$ are retained.
\end{theorem}
\begin{proof}
Put $C=A_{HH}$ and $B=A_{H,H^c}$.  If $s_3(C)\ge a$,
condition on the complementary original coordinates and apply
Corollary~\ref{cor:finite-core}.  This gives the first term.
If $s_3(C)<a$ and $s_2(C)\ge b$,
Lemma~\ref{supp:pair-extraction} gives a physical pair $K$ with
$s_{\min}(K)>\kappa$, hence $|\det K|>\kappa^2$.
A definite pair gives $\pi/\kappa$ by Lemma~\ref{geom:definite}.
An indefinite pair gives the estimate in Theorem~\ref{geom:fullpair}.
The function
\[
 x\longmapsto\frac{9.2+\log_+(10000x/s)}x
\]
is decreasing on $(0,\infty)$: on its logarithmic branch its
derivative has numerator $1-9.2-\log(10000x/s)<0$, and
elsewhere this is immediate.  Substitute $x=\sqrt{|\det K|}
>\kappa$; the second term in \eqref{eq:supp-six-carrier}
also exceeds $\pi/\kappa$.

The remaining case has $s_2(C)<b$.
Lemma~\ref{supp:spill} with $k=6,r=1$ gives
$s_5(B)>\sqrt{g^2-b^2}$.
The five-volume case of Corollary~\ref{four:d-cases} gives
\[
 \|p_Q\|_\infty\le5\sqrt{2e}\,
       e_5(BB^{\mathsf T})^{-1/10}
 \le5\sqrt{2e}/s_5(B),
\]
which proves the third term.  These three alternatives are exhaustive.
\end{proof}

The carrier hypothesis itself implies $s_6(A)\ge g$, so one may
always take $s=g$ if no stronger global spectral lower bound is
available.  In particular, this theorem does not extract six directions
from the three-direction premise of Conjecture~\ref{conj:main}.

\subsection{Stability under additional spectrum}

\begin{proposition}[Localized anticommutator bound]
\label{supp:anticommutator}
Let $A=A_0+R$ be symmetric and define
\[
 \nu_H=\max\{0,-\lambda_{\min}((A_0R+RA_0)_{HH})\}.
\]
For every $j\le|H|$,
\[
 \lambda_j((A^2)_{HH})\ge
       \lambda_j((A_0^2)_{HH})-\nu_H.
\]
If $A_0R+RA_0=0$, all these inequalities hold with zero loss.
This includes disjoint spectral packets, for which $A_0R=0$.
\end{proposition}
\begin{proof}
The identity $A^2=A_0^2+A_0R+RA_0+R^2$, together with
$R^2\ge0$, gives $(A^2)_{HH}\ge(A_0^2)_{HH}-\nu_HI$.
Apply monotonicity of ordered eigenvalues under this order.
The zero-loss assertions follow from the same identity.
\end{proof}

\begin{theorem}[Low-rank additions with an operator-norm remainder]
\label{supp:rank-reserve}
Let $A=A_0+S+R$ be symmetric, $\|R\|\le\varepsilon$, and
$r_H=\operatorname{rank}(SP_H)$.  If $k+r_H\le|H|$, then
\begin{equation}\label{eq:supp-rank-reserve}
 \sqrt{\lambda_k((A^2)_{HH})}\ge
 \left(\sqrt{\lambda_{k+r_H}((A_0^2)_{HH})}
                                  -\varepsilon\right)_+.
\end{equation}
No restriction is imposed on the signs or magnitudes of the
nonzero eigenvalues of $S$.
\end{theorem}
\begin{proof}
For any rectangular matrices $T,E$ with $\operatorname{rank}E\le r$,
the singular-value variational principle gives
\begin{equation}\label{eq:supp-rank-interlace}
 s_k(T+E)\ge s_{k+r}(T).
\end{equation}
Indeed, take a $(k+r)$-dimensional subspace on which
$|Tx|\ge s_{k+r}(T)|x|$ and intersect it with $\ker E$.
The intersection has dimension at least $k$, and $T+E=T$ there.
Apply \eqref{eq:supp-rank-interlace} to $T=A_0P_H$, $E=SP_H$,
then use $\|RP_H\|\le\varepsilon$.  Finally,
\[
 s_j(MP_H)^2=\lambda_j((M^2)_{HH})
\]
for symmetric $M$, which proves \eqref{eq:supp-rank-reserve}.
\end{proof}

\begin{remark}[The amount of spectral reserve]
A different argument notes that
\[
 \operatorname{range}((A_0+S)^2-A_0^2)
 \subset \operatorname{range}(S)+A_0\operatorname{range}(S).
\]
Rank interlacing of the squares then gives a valid but weaker loss
of at most $2\operatorname{rank}S$.  The fact that this difference
can have rank $2\operatorname{rank}S$ does not make that spectral
reserve loss necessary.  Theorem~\ref{supp:rank-reserve} uses the
rectangular operators before squaring and loses at most
$\operatorname{rank}S$, and often only $r_H$.
\end{remark}

\begin{corollary}[An adaptive split of a signed addition]
\label{supp:adaptive-rank}
Let $E=E^{\mathsf T}$, let $S_h$ contain its eigenmodes with
absolute eigenvalue greater than $h\ge0$, and put
$r_H(h)=\operatorname{rank}(S_hP_H)$.  If $\|T\|\le\delta$
and $A=A_0+E+T$, then, whenever $k+r_H(h)\le|H|$,
\[
 \sqrt{\lambda_k((A^2)_{HH})}\ge
 \left(\sqrt{\lambda_{k+r_H(h)}((A_0^2)_{HH})}
                                      -h-\delta\right)_+.
\]
The threshold $h$ may be chosen to maximize the bound.
\end{corollary}
\begin{proof}
The remaining part $E-S_h$ has operator norm at most $h$.
Apply Theorem~\ref{supp:rank-reserve} with remainder
$E-S_h+T$.
\end{proof}

\begin{proposition}[Separate positive and negative reserves]
\label{supp:signed-reserve}
Let $C=C_0+P-N+T$ be symmetric $m$ by $m$, where $P,N\ge0$
and $\|T\|\le\delta$.  Write $p_j=\lambda_j(P)$ and
$n_j=\lambda_j(N)$ in decreasing order, and set
$p_{m+1}=n_{m+1}=0$.  For $1\le u,v\le m$,
\begin{align*}
 \lambda_u(C)&\ge
 \max_{0\le r\le m-u}
    \{\lambda_{u+r}(C_0)-n_{r+1}\}-\delta,\\
 \lambda_v(-C)&\ge
 \max_{0\le r\le m-v}
    \{\lambda_{v+r}(-C_0)-p_{r+1}\}-\delta.
\end{align*}
If the active bounds are positive and supply a total of at least
three positive and negative directions, their minimum is a lower
bound for $s_3(C)$.
\end{proposition}
\begin{proof}
Write $N=N_r+N_{\mathrm{tail}}$, where $N_r$ has rank at most $r$
and $\|N_{\mathrm{tail}}\|=n_{r+1}$.  The favorable term $P\ge0$
can only increase ordered eigenvalues, so
\[
 \lambda_{u+r}(C_0+P+T-N_{\mathrm{tail}})
 \ge\lambda_{u+r}(C_0)-\delta-n_{r+1}.
\]
Subtracting $N_r$ loses at most $r$ indices by rank interlacing,
which proves the first inequality after maximizing in $r$.
Apply the same argument to $-C$ for the second.  Strictly positive
positive and negative spectral directions are disjoint, proving
the final assertion.
\end{proof}

\subsection{A complementary Schur block protects three directions}

\begin{theorem}[Complementary Schur protection]
\label{supp:schur-protection}
On an original coordinate block $H$ of size $m$, suppose
\[
 A_{HH}=B+U\Theta U^{\mathsf T}+T,
 \qquad B=B^{\mathsf T},\quad\Theta=\Theta^{\mathsf T},
 \quad\|T\|\le\delta.
\]
Let the columns of $W$ form an orthonormal basis of
$\operatorname{range}U$, and let the columns of $Z$ form an
orthonormal basis of its orthogonal complement.
Suppose $D=Z^{\mathsf T}BZ$ is invertible and $s_3(D)\ge\mu>0$.
Put
\[
 R=(W^{\mathsf T}BZ)D^{-1},\quad\rho=\|R\|,\quad
 \ell=\frac{\sqrt{\rho^2+4}-\rho}{2}.
\]
If $\delta<\ell^2\mu$, then for every $\Theta$,
\[
 \|p_Q\|_\infty\le
 \frac{2\sqrt{6\pi m}}{\ell^2\mu-\delta}.
\]
\end{theorem}
\begin{proof}
The orthogonal basis $[W\ Z]$ is used for matrix algebra only.
In this basis the block without $T$ equals
\[
 L\begin{pmatrix}S&0\\0&D\end{pmatrix}L^{\mathsf T},
 \quad L=\begin{pmatrix}I&R\\0&I\end{pmatrix},
 \quad
 S=W^{\mathsf T}(B+U\Theta U^{\mathsf T})W-RDR^{\mathsf T}.
\]
No invertibility assumption on $S$ is needed.  A singular-value
decomposition of $R$ decomposes $L$ into two-dimensional shears
$\bigl(\begin{smallmatrix}1&r\\0&1\end{smallmatrix}\bigr)$
and identity blocks.  The smaller singular value of such a shear
is $(\sqrt{r^2+4}-r)/2$, decreasing in $r\ge0$, so
$s_{\min}(L)=\ell$.  Multiplication on each side by an invertible
matrix reduces every singular value by at most its smallest singular
value.  Consequently
\[
 s_3(B+U\Theta U^{\mathsf T})
 \ge\ell^2s_3(\operatorname{diag}(S,D))\ge\ell^2\mu.
\]
The perturbation $T$ costs at most $\delta$.  Condition outside $H$
and apply Corollary~\ref{cor:finite-core} to finish.
\end{proof}

\begin{corollary}[Many large eigenvalues of one sign]
\label{supp:codimension-one}
If $A$ has at least $n-1$ eigenvalues at least $s>0$, or at least
$n-1$ eigenvalues at most $-s$, and $n\ge4$, then
\[
 \|p_Q\|_\infty\le\frac{2\sqrt{24\pi}}s.
\]
The remaining eigenvalue may have any sign and magnitude.
\end{corollary}
\begin{proof}
Every principal four-coordinate block has at least three eigenvalues
of the indicated sign and magnitude by interlacing.  Apply
Corollary~\ref{cor:finite-core} with $m=4$.
\end{proof}

These results give bounds under additional spectral hypotheses.
They do not establish that $s_3(A)\ge1$ alone guarantees a bounded
coordinate block, linear bank, or complementary Schur block satisfying
their respective hypotheses.

\subsection{Limitations of density interpolation and joint caps}

\begin{proposition}[An $L^3$ bound and a square-root interval bound]
\label{supp:interpolation-obstruction}
There are probability densities $p_N$, $N\ge2$, such that
\[
 \|p_N\|_3\le2^{1/3},\qquad
 \sup_z\int_{z-\varepsilon}^{z+\varepsilon}p_N(x)\,dx
        \le3\sqrt\varepsilon\quad(\varepsilon>0),
\]
while $\|p_N\|_\infty=N$.
\end{proposition}
\begin{proof}
Set
\[
 p_N=N\mathbf1_{[0,N^{-3}]}
                +(1-N^{-2})\mathbf1_{[2,3]}.
\]
Its mass is one, its essential supremum is $N$, and
$\int p_N^3=1+(1-N^{-2})^3\le2$.
An interval of length $2\varepsilon$ receives at most
$\min\{N^{-2},2N\varepsilon\}$ from the spike.  Dividing by
$\sqrt\varepsilon$, the maximum over $\varepsilon>0$ occurs
at $\varepsilon=(2N^3)^{-1}$ and equals
$\sqrt2N^{-1/2}\le1$.  If $0<\varepsilon\le1$, the other
part has mass at most $\min\{1,2\varepsilon\}\le2\sqrt\varepsilon$.
For $\varepsilon\ge1$, the total mass is already at most
$1\le3\sqrt\varepsilon$.  This proves all assertions.
\end{proof}

Thus a uniform $L^3$ bound and a uniform square-root interval bound
do not imply a uniform density bound.  An additional hypothesis is
needed to combine these estimates for quadratic forms.

\begin{proposition}[A moving quadratic ridge under a joint density cap]
\label{supp:moving-ridge}
For every $0<\varepsilon<1$ there are a uniform $X$ on $[0,1]$
and a triple $(Y,Z,W)$ independent of $X$, with joint probability density
bounded by one such that
\[
 T=4XY+Z^2-YW,\qquad \|p_T\|_\infty\ge\frac1{9\varepsilon}.
\]
The quadratic matrix has spectrum
$\{1,0,\sqrt{17}/2,-\sqrt{17}/2\}$.  The triple is not required
to have independent coordinates.
\end{proposition}
\begin{proof}
For $L>0$, define
\[
 S_{\varepsilon,L}=
 \{(y,z,w):\varepsilon\le y\le2\varepsilon,
       \ L\le w\le2L,
       \ \sqrt{yw}\le z\le\sqrt{yw+\varepsilon}\}.
\]
Its volume
\[
 V_\varepsilon(L)=\int_\varepsilon^{2\varepsilon}
       \int_L^{2L}
       (\sqrt{yw+\varepsilon}-\sqrt{yw})\,dw\,dy
\]
is positive and continuous.  It tends to zero as $L\downarrow0$,
because it is at most $\varepsilon^{3/2}L$.
For $L\ge1$, rationalization and $yw\le4\varepsilon L$ give
\[
 V_\varepsilon(L)\ge
    \frac{\varepsilon^{3/2}\sqrt L}{2\sqrt5},
\]
so it tends to infinity as $L\to\infty$.  Choose $L$ for which
$V_\varepsilon(L)=1$ and let the triple have density
$\mathbf1_{S_{\varepsilon,L}}$.
On this set, $0\le Z^2-YW\le\varepsilon$, while
$0\le4XY\le8\varepsilon$.  Thus $T\in[0,9\varepsilon]$
almost surely.  Given the triple, $4XY$ is uniform on an interval
of positive length; consequently $T$ has a density.  A probability
density supported on an interval of length $9\varepsilon$ has
essential supremum at least its reciprocal.

In coordinates $(X,Y,Z,W)$ the matrix is
\[
 A_*=\begin{pmatrix}
 0&2&0&0\\2&0&0&-1/2\\0&0&1&0\\0&-1/2&0&0
 \end{pmatrix}.
\]
The $(X,Y,W)$ block has eigenvalues
$0,\pm\sqrt{2^2+(1/2)^2}$, and $Z$ contributes eigenvalue one.
This proves the stated spectrum.
\end{proof}

Proposition~\ref{supp:moving-ridge} concerns a weaker independence
assumption than Conjecture~\ref{conj:main}.  It shows that keeping
one independent hub and a bounded joint density for the remaining
three statistics does not preserve enough information for the
desired conclusion.  It is not a counterexample to that conjecture.
